\Article{17}
The following equation thus expresses the condition that the combined forces do not turn the movable needle about the vertical axis.
\[
0=-mT\sin u+fR^{-(n+1)}+f'R^{-(n+2)}+f''R^{-(n+3)}+\text{etc.}-\theta(u-N)
\]

It is easy to ensure that the value of $N$, if not exactly $=0$, is at least very small. In the experiments in question, $u$ also remains within narrow \SourceLeaf{109}limits. We can therefore replace the term $\theta(u-N)$ by $\theta\sin(u-N)$ without fear of appreciable error, especially since $\frac{\theta}{mT}$ is a very small fraction. Let $u^0$ be the value of $u$ corresponding to equilibrium of the first needle in the absence of the second, or equivalently let
\[
mT\sin u^0+\theta\sin(u^0-N)=0
\]
from which it readily follows that
\[
mT\sin u+\theta\sin(u-N)=(mT\cos u^0+\theta\cos(u^0-N))\sin(u-u^0)
\]
Here we can safely replace the first factor by $mT+\theta$. Our equation thus becomes
\[
(mT+\theta)\sin(u-u^0)=fR^{-(n+1)}+f'R^{-(n+2)}+f''R^{-(n+3)}+\text{etc.}
\]
If we retain only the first term $fR^{-(n+1)}$ here, the solution is readily obtained. We have
\[
\begin{aligned}
&\operatorname{tang}(u-u^0)\\
&=\frac{\begin{gathered}mM(n\cos(\psi-U)\sin(\psi-u^0)\\+\sin(\psi-U)\cos(\psi-u^0))R^{-(n+1)}\end{gathered}}{\begin{gathered}mT+\theta+mM(n\cos(\psi-U)\cos(\psi-u^0)\\-\sin(\psi-U)\sin(\psi-u^0))R^{-(n+1)}\end{gathered}}
\end{aligned}
\]
We are just as well justified in dropping the part of the denominator that contains the factor $R^{-(n+1)}$, or in setting
\[
\begin{aligned}
\operatorname{tang}(u-u^0)&=\frac{mM}{mT+\theta}(n\cos(\psi-U)\sin(\psi-u^0)\\
&\qquad+\sin(\psi-U)\cos(\psi-u^0))R^{-(n+1)}\\
&=FR^{-(n+1)}
\end{aligned}
\]

If we wish to take the further terms into account, it is clear that $\operatorname{tang}(u-u^0)$ expands into a series of the form
\[
\operatorname{tang}(u-u^0)=FR^{-(n+1)}+F'R^{-(n+2)}+F''R^{-(n+3)}+\text{etc.}
\]
A moment's thought shows that the coefficients $F,F',F''$ and so on, up to and including the coefficient of the power $R^{-(2n+1)}$, arise respectively from
\[
\frac{f}{mT+\theta},\quad\frac{f'}{mT+\theta},\quad\frac{f''}{mT+\theta}\quad\text{etc.}
\]
by replacing $u$ with $u^0$. From the next term onward, new contributions appear, which we need not pursue more closely for our purpose. Moreover, it is clear that $u-u^0$ will expand into a series of similar form, which agrees even as far as the power $R^{-(3n+2}$ with the series for $\operatorname{tang}(u-u^0)$.

